\section{Literature Review}\label{sec:lit}

\subsection{Machine Learning for Crop-Yield Prediction}

Statistical crop--climate models, in which yield is regressed on seasonal weather summaries,
have long been used as interpretable baselines \cite{lobell2010}. In a systematic review of 50
machine-learning studies, temperature, rainfall and soil type were found to be the most frequently
used features, and neural networks, random forests and gradient boosting the most frequently used
models \cite{vanklompenburg2020}. Random forests were shown to outperform multiple linear regression
for global and regional yield prediction \cite{jeong2016}, and deep networks were applied to large
maize datasets \cite{khaki2019}. For India, district yields from the Directorate of Economics and
Statistics have been combined with NASA POWER weather. In a study of pearl millet in nine districts
of Rajasthan from 1997 to 2019, no single model was best in every district, and XGBoost and
generalised linear models fitted the training years almost perfectly but often failed on the test
years \cite{alsaber2025}. In a benchmark of nine model classes on 275 wheat districts from 2001
to 2023, a hybrid CNN--BiLSTM--attention model reached an $R^2$ of 0.795, but a persistence
forecast that repeats the previous year's yield was almost as accurate on average
\cite{manogna2026}. Both studies tested on later years of the same districts and assumed complete
weather inputs.

\subsection{Missing Data in Agricultural Prediction}

Missing values are classified as missing completely at random (MCAR), missing at random, or
missing not at random \cite{little2019}. Common remedies include replacement with a mean or median,
$k$-nearest-neighbour (KNN) imputation \cite{troyanskaya2001}, chained-equation (iterative)
imputation \cite{vanbuuren2011} and random-forest imputation \cite{stekhoven2012}, while
gradient-boosted trees such as XGBoost can route missing values through learned default branches
\cite{chen2016}. For Indian wheat data from 370 districts, Saini and Nagpal \cite{saini2024} found
that simple mean replacement performed best among the statistical methods ($R^2 = 0.83$) and
MissForest and MICE among the learned methods ($R^2 = 0.80$). For sunflower in Romania, Călin
\emph{et al.} \cite{calin2023} raised $R^2$ from 0.723, obtained when incomplete rows were removed,
to 0.938 with random-forest imputation and outlier removal. In both studies, the gaps that
happened to be present in the data were used, so the rate and the pattern of missingness were not
controlled.

\subsection{Spatial Validation and Distribution Shift}

When observations close in space are split between training and test sets, random
cross-validation overestimates the accuracy that can be expected at new locations
\cite{roberts2017,meyer2019}, and spatial validation of large-scale ecological maps has revealed
much poorer performance than random validation \cite{ploton2020}. Such distribution shifts between
training and deployment domains have been formalised in benchmarks such as WILDS \cite{koh2021}.
For crop yield, Chakravarty \cite{chakravarty2026} showed on more than 1,200 United States counties
that leave-one-region-out evaluation causes large performance gaps between agro-ecological regions,
with negative $R^2$ for some models. In Germany, models that performed well on spatial test splits
lost accuracy on independent years, and their SHAP importances remained plausible even when they
did not generalise \cite{baatz2025}. In Sub-Saharan Africa, leave-one-country-out evaluation gave
negative $R^2$ for every feature set, whereas within-country random cross-validation gave moderate
values \cite{adjei2026}.

\subsection{Research Gap}

The closest studies are summarised in Table~\ref{tab:related}. Indian district studies have assumed
complete weather and tested only on familiar districts; imputation studies have not controlled the
missingness mechanism or tested transfer to new regions; and studies of geographic shift have not
considered missing inputs and have not used Indian data. To the best of the authors' knowledge, a
controlled evaluation of missing weather data and geographic shift \emph{together}, which also
measures how much random-split evaluation overstates accuracy, has not been reported for Indian
district-level yield prediction. This gap is addressed in the present study.

\begin{table}[!t]
\caption{Summary of Closely Related Studies}\label{tab:related}
\centering\footnotesize
\setlength{\tabcolsep}{3pt}
\begin{tabular}{p{1.55cm}p{1.6cm}p{1.35cm}ccc}
\toprule
Study & Data & Models & Ctrl.\ miss. & Unseen region & Joint \\
\midrule
Alsaber \emph{et al.} \cite{alsaber2025} & Rajasthan, bajra, 9 distr. & GLM, RF, XGB, SVR & No & No & No \\
Manogna R.~L. \emph{et al.} \cite{manogna2026} & India, wheat, 275 distr. & RF, boosting, deep & No & No & No \\
Saini \& Nagpal \cite{saini2024} & India, wheat, 370 distr. & 8 imputers & No & No & No \\
Călin \emph{et al.} \cite{calin2023} & Romania, sunflower & 9 imputers, boosting & No & No & No \\
Chakravarty \cite{chakravarty2026} & USA, corn, soybean & GNN-RNN, ViT & No & Yes & No \\
Baatz \cite{baatz2025} & Germany, NUTS-3 & RF, XGB, LSTM & No & Yes & No \\
Adjei \cite{adjei2026} & Africa, maize & Embeddings & No & Yes & No \\
\textbf{This work} & Gujarat, 3 crops, 32 distr. & Ridge, RF, XGB & \textbf{Yes} & \textbf{Yes} & \textbf{Yes} \\
\bottomrule
\end{tabular}

\vspace{2pt}{\scriptsize Ctrl.\ miss.: missingness simulated at controlled rates and patterns;
Joint: both conditions evaluated together.}
\end{table}
